The second variety of interleaved methods involves navigating the state and control spaces using forward control propagation. Using a random state and action, these methods determine the resultant state, constructing a tree of dynamically feasible edges \cite{pasricha2022pokerrt}.
SyCLoP and KPIECE leverage state space discretizations and guide exploration in low-coverage areas \cite{plaku2010motion,sucan2011sampling}.
SST optimizes path quality via pruning and heuristic biasing \cite{li2016asymptotically}.
GBRRT and GABRRT combine multiple propagation methods and backward tree costs as heuristic values for bidirectional kinodynamic search \cite{nayak2022bidirectional}. Other techniques involve pre-mapping the space of valid motions and conducting a search over this state and control space discretization \cite{shome2021asymptotically,bordalba2018randomized}.
Existing approaches that exploit lazy strategies for collision checking involve learning-based methods \cite{yu2021reducing} and evaluating paths on an as-needed basis \cite{kavraki2000path,mandalika2019generalized,hauser2015lazy}. We take an empirical approach by computing a collision probability for each kinodynamic action in our state and control space discretization.

Overall, interleaved approaches are easy to design since they leverage general-purpose physics simulations, explore both state and control spaces uniformly, and easily adapt to environmental interactions and multi-object systems. However, they may face slow convergence owing to costly simulation and collision checking methods. Next, we outline how our contribution mitigates these issues, paving the way for fast computation of asymptotically optimal solutions.
